\documentclass[10pt,a4paper]{amsart}
\usepackage[margin=19mm]{geometry}
\usepackage{amsmath,amssymb,mathtools}
\usepackage[scr=rsfs,cal=euler]{mathalfa}
\usepackage{xfrac}
\usepackage[unicode,hidelinks,bookmarks=false]{hyperref}
\newcommand{\ie}{i.e.\ }
\newcommand{\cf}{cf.\ }
\newcommand{\resp}{respectively\ }
\newcommand{\Subsection}{\subsection}
\providecommand{\nobreakdash}{\nobreak\hbox{-}}
\setlength{\emergencystretch}{2em}
\multlinegap0pt
\usepackage{cleveref}

\theoremstyle{plain}
\newtheorem{theorem}{Theorem}[section]
\theoremstyle{remark}
\newtheorem{remark}{Remark}[section]
\newtheorem{conclusion}[remark]{Conclusion}
\newtheorem{condition}{Condition}[section]
\theoremstyle{definition}
\newtheorem{definition}{Definition}[section]
\newtheorem{example}{Example}[section]
\theoremstyle{plain}
\newtheorem{assumption}{Assumption}[section]
\newtheorem{problem}{Problem}
\newtheorem{lemma}{Lemma}
\newtheorem{corollary}{Corollary}
\newtheorem{proposition}{Proposition}
\title[Null Controllability of a Stochastic Parabolic Equation with a Space-Dependent Analytic Noise Coefficient]{Null Controllability of a Stochastic Parabolic Equation with a Space-Dependent Analytic Noise Coefficient}
\author{Qi L\"u, Yu Wang}
\date{Source: arXiv:2609.28937v1 (CC BY 4.0)}
\begin{document}
\maketitle

\noindent\emph{Source and licence.} The delivered work is the article shown in the title above, version arXiv:2609.28937v1 (CC BY 4.0), distributed under the Creative Commons Attribution 4.0 International licence, as recorded in the arXiv licence metadata for this exact version and shown on the abstract page saved with this item. The full licence notice, which carries the licence URL and the address of that abstract page, is the bundled file \texttt{licenses/arxiv-2609.28937-CC-BY-4.0.txt}.
\par\smallskip

\noindent\textit{Editorial scope.} Counted unit: the article's Proposition 3.2 (source0.tex lines 669--683, source label \texttt{prop:hilbert-propagation}): ``Periodic Hilbert space-valued propagation Let (H ) be a separable complex Hilbert space, let (r 0 ), and let (E T n ) be a measurable set of positive measure. Then there exist constants (C 0 ) and ( (''. It is delivered together with the article's complete original proof (source0.tex lines 686--754, one proof environment printed later in the article, detached from the statement by the article's own arrangement, 69 lines), copied line for line from the article's own text. Required context retained uncounted: thm:scalar-propagation (lines 648--663). Editorial formatting only: an AMS \texttt{amsart} preamble that restates the article's own macro definitions and its own theorem declarations, and the theorem environments the retained lines use; the packages cleveref are loaded to match the article's own setup; the article's own class file is neither shipped nor input; the retained environments are renumbered sequentially in order of appearance, whereas the article prints them with its own numbering; the article's equation numbering is not reproduced and every cross-reference inside the retained text resolves to the wrapper's numbering. No citation occurs inside any retained line, so the wrapper carries no bibliography; All retained bodies are exact copies of the bundled \texttt{sources/211626/source0.tex}, so no omission receipt applies. No asset-inclusion command occurs inside any retained span and the package ships no image asset.


\section*{Retained context: thm:scalar-propagation (source0.tex lines 648--663)}

\begin{lemma}[Multidimensional analytic propagation]
\label{thm:scalar-propagation}
Let \(\ell>0\), let \(E\subset B_\ell(x_0)\subset\mathbb R^n\) be measurable with positive Lebesgue measure \(|E|\), and let \(f:B_{2\ell}(x_0)\to\mathbb R\) be real-analytic. Assume that, for some \(\mathcal M>0\) and \(0<\varrho\le1\),
\[
|\partial^\alpha f(x)|
\le \mathcal M\frac{|\alpha|!}{(\varrho\ell)^{|\alpha|}},
\qquad x\in B_{2\ell}(x_0),\quad \alpha\in\mathbb N_0^n.
\]
Then there exist constants \(N\ge1\) and \(\vartheta\in(0,1)\), depending only on \(n\), \(\varrho\), and \(|E|/|B_\ell(x_0)|\), such that
\[
\|f\|_{L^\infty(B_\ell(x_0))}
\le N
\left(\frac1{|E|}\int_E|f(x)|\,dx\right)^\vartheta
\mathcal M^{1-\vartheta}.
\]
\end{lemma}


\section{The counted unit: Proposition 3.2 and its complete original proof}

\begin{proposition}[Periodic Hilbert space-valued propagation]
\label{prop:hilbert-propagation}
Let \(H\) be a separable complex Hilbert space, let \(r>0\), and let \(E\subset\mathbb T^n\) be a measurable set of positive measure. Then there exist constants \(C>0\) and \(\theta\in(0,1)\), depending only on \(n\), \(r\), and \(|E|\), such that the following holds. If \(f:S_r^n\to H\) is holomorphic, \(2\pi\)-periodic in each coordinate, and satisfies
\[
\sup_{\zeta\in S_r^n}\|f(\zeta)\|_H\le\mathcal M
\]
for some \(\mathcal M\ge0\), then
\begin{equation}\label{eq:hilbert-propagation}
\|f\|_{L^2(\mathbb T^n;H)}
\le C
\|f\|_{L^2(E;H)}^\theta
\mathcal M^{1-\theta}.
\end{equation}
In particular, the constants \(C\) and \(\theta\) do not depend on \(H\).
\end{proposition}

\begin{proof}
If \(\mathcal M=0\) or \(H=\{0\}\), the conclusion is immediate. Hence we may assume that \(\mathcal M>0\) and \(H\ne\{0\}\). The proof is divided into two steps.

\smallskip

\textbf{Step 1.} We first apply the scalar propagation estimate.
Regard \(E\) as a subset of \(Q=[-\pi,\pi)^n\), and set \(\varrho=\min\{1,r/(4\pi\sqrt n)\}\). Then
\(E\subset Q\subset B_{2\pi\sqrt n}(0)\) and \(2\pi\sqrt n\varrho \le r/2\).
For each unit vector \(h\in H\) and each \(\gamma\in\mathbb R\), define
\begin{align*}
g_{h,\gamma}(x)
=\operatorname{Re}\bigl(e^{-\mathrm{i}\gamma}\langle f(x),h\rangle_H\bigr),
\qquad x\in\mathbb R^n.
\end{align*}
Since the inner product is linear in its first variable, the function
\(\zeta\mapsto\langle f(\zeta),h\rangle_H\) is holomorphic on \(S_r^n\) and bounded by \(\mathcal M\).
Every polydisc of radius \(r/2\) centered at a real point is contained in \(S_r^n\).
Thus Cauchy's inequality yields, for \(x\in B_{4\pi\sqrt n}(0)\),
\begin{align*}
|\partial^\alpha g_{h,\gamma}(x)|
&\leq |\langle \partial^\alpha f(x),h\rangle_H|
\leq \mathcal M\alpha!\biggl(\frac2r\biggr)^{|\alpha|}
\leq \mathcal M\frac{|\alpha|!}{(2\pi\sqrt n\varrho )^{|\alpha|}}.
\end{align*}
Here we used \(\alpha!\le|\alpha|!\) and \(2\pi\sqrt n\varrho \le r/2\).
Therefore \Cref{thm:scalar-propagation} applies and gives
\begin{align*}
\sup_{x\in B_{2\pi\sqrt n}(0)}|g_{h,\gamma}(x)|
\leq C\biggl(\frac1{|E|}\int_E|g_{h,\gamma}(x)|\,dx\biggr)^\theta
\mathcal M^{1-\theta}.
\end{align*}
For the observation term, we have
\begin{align*}
\frac1{|E|}\int_E|g_{h,\gamma}(x)|\,dx
&\leq\frac1{|E|}\int_E\|f(x)\|_H\,dx
\leq |E|^{-1/2}\|f\|_{L^2(E;H)}.
\end{align*}
Consequently,
\begin{align*}
\sup_{x\in B_{2\pi\sqrt n}(0)}|g_{h,\gamma}(x)|
\leq C|E|^{-\theta/2}
\|f\|_{L^2(E;H)}^\theta\mathcal M^{1-\theta},
\end{align*}
where \(C\) and \(\theta\) depend only on \(n\), \(r\), and \(|E|\), and are independent of \(h\), \(\gamma\), and \(H\).

\smallskip

\textbf{Step 2.} We now recover the Hilbert-space norm.
For every \(x\in Q\), the dual characterization of the norm yields
\begin{align*}
\|f(x)\|_H
&=\sup_{\|h\|_H=1}|\langle f(x),h\rangle_H|
=\sup_{\|h\|_H=1,\,\gamma\in\mathbb R}
\bigl|\operatorname{Re}\bigl(e^{-\mathrm{i}\gamma}\langle f(x),h\rangle_H\bigr)\bigr|.
\end{align*}
Taking the supremum over \(h\) and \(\gamma\) in the estimate from Step~1, we obtain
\begin{align*}
\sup_{x\in Q}\|f(x)\|_H
\leq C|E|^{-\theta/2}
\|f\|_{L^2(E;H)}^\theta\mathcal M^{1-\theta}.
\end{align*}
Since \(f\) is periodic, integrating over \(Q\) gives
\begin{align*}
\|f\|_{L^2(\mathbb T^n;H)}
&=\biggl(\int_Q\|f(x)\|_H^2\,dx\biggr)^{1/2}\leq |Q|^{1/2}\sup_{x\in Q}\|f(x)\|_H
\leq C\|f\|_{L^2(E;H)}^\theta\mathcal M^{1-\theta}.
\end{align*}
This proves \eqref{eq:hilbert-propagation}, with constants independent of \(H\).
\end{proof}

\end{document}
